% !TeX root = ../navier-stokes-manuscript.tex
% ============================================================
% 2.1 Critical Height at a Finite Endpoint
% ============================================================

\subsection{Critical Height at a Finite Endpoint}\label{sub:CriticalHeightAtAFiniteEndpoint}

Return to the one maximal solution generated by the rapidly decaying datum
\(u_0\) at fixed \(\nu>0\), and suppose its last time \(T_*\) is finite.  The
field is smooth before that time.  Its threat must therefore be a change of
scale along the same history, not an already singular initial state.

Kinetic energy remains bounded, yet a structure’s contribution to that energy
changes when its width changes.  For \(\lambda>0\), set
\[
  u_\lambda(x,t)
  :=
  \lambda u(\lambda x,\lambda^2t),
  \qquad
  p_\lambda(x,t)
  :=
  \lambda^2p(\lambda x,\lambda^2t)
\]
These fields solve Navier--Stokes with the same viscosity on the rescaled time
interval.  Space contracts by \(\lambda^{-1}\), time by \(\lambda^{-2}\), and
velocity grows by \(\lambda\); one time derivative and the viscous two-space-
derivative term consequently transform together.

Let \(\Lambda=(-\Delta)^{1/2}\) and measure height by
\[
  \norm{v}_{\dot H^s}^2
  :=
  \int_{\R^3}
  |\xi|^{2s}|\widehat v(\xi)|^2
  \,d\xi.
\]
\begin{proposition}[Critical Sobolev scaling]\label[proposition]{prop:CriticalSobolevScaling}
Whenever the homogeneous norm is finite, every \(s\in\R\) and \(\lambda>0\)
satisfy
\[
  \norm{u_\lambda(t)}_{\dot H^s}^2
  =
  \lambda^{2s-1}
  \norm{u(\lambda^2t)}_{\dot H^s}^2.
\]
The only scale-invariant member is \(s=\tfrac12\).
\end{proposition}
\begin{proof}
The rescaled Fourier transform is
\[
  \widehat{u_\lambda}(\xi,t)
  =
  \lambda^{-2}
  \widehat u(\xi/\lambda,\lambda^2t).
\]
Substitution into the norm, followed by \(\eta=\xi/\lambda\), gives
\begin{align*}
  \norm{u_\lambda(t)}_{\dot H^s}^2
  &={}
  \int_{\R^3}
  |\xi|^{2s}\lambda^{-4}
  \left|\widehat u(\xi/\lambda,\lambda^2t)\right|^2
  \,d\xi
  \\
  &={}
  \lambda^{2s-1}
  \int_{\R^3}
  |\eta|^{2s}
  \left|\widehat u(\eta,\lambda^2t)\right|^2
  \,d\eta.
\end{align*}
The prefactor is one exactly when \(2s-1=0\), proving the claim.
\end{proof}
Compression therefore tilts every Sobolev height except \(1/2\).  Name the
quantity that remains unchanged:
\[
  H(t)
  :=
  \frac12\norm{u(t)}_{\dot H^{1/2}}^2
  =
  \frac12\norm{\Lambda^{1/2}u(t)}_2^2.
\]
An exact rescaling can now narrow and accelerate a smooth configuration without
changing \(H\).

At energy height, the same narrowing has the opposite weight:
\[
  \norm{u_\lambda(t)}_2^2
  =
  \lambda^{-1}
  \norm{u(\lambda^2t)}_2^2.
\]
A finer structure can therefore retain its critical height while carrying less
kinetic energy.  The energy ceiling does not itself prevent the active width
from continuing to fall.

The parabolic durations in \cref{sub:HistoricalEstimates} allow successively
finer scales to fit arithmetically before a finite endpoint.  At critical
height, those rescaled copies can all have the same \(H\) while their energy
falls.  They are still only copies: scaling has not made the actual solution
visit them.  Such a passage would have to be generated by transport, pressure,
incompressibility, and viscosity acting together on the preceding state.

The scale is now fixed.  A finite endpoint turns the remaining issue into
response: if one row of the actual field escapes, what must the next time row do
before the available time runs out?
